\subsection*{Intended Learning Outcomes - Lecture 10}

The student should be able to:
\begin{itemize}
\item know the definitions of indirect/direct causal relations, confounder and confounding according to simple SCMs;
\item know how to detect the presence of such relations from Markov kernels induced by simple SCMs;
\item be able to express the notion of ``no confounding'' in terms of potential outcomes;
\item know how to deal formally with a situation in which only part of the variables in an SCM are observed, whereas the others are latent;
\item be able to formally model an RCT with SCMs in various ways, and to prove in the SCM formalism the main properties of RCTs;
\item know the definitions of faithfulness;
\item understand different ways of how faithfulness can occur;
\item understand the LCD algorithm;
\item be able to prove the correctness of the LCD algorithm;
\item understand the Y-structures approach to causal discovery;
\item be able to prove the correctness of the Y-structures approach to causal discovery.
\end{itemize}
